% Part: lambda-calculus
% Chapter: syntax
% Section: eta

\documentclass[../../../include/open-logic-section]{subfiles}

\begin{document}

\olfileid{lam}{syn}{eta}
\olsection{تبدیل~$\eta$}

رابطهٔ دیگری نیز روی ترم‌های~$\lambda$ وجود دارد. در~\olref[fv]{sec}
از مثال~$\lambd[x][(fx)]$ استفاده کردیم؛ این ترم آرگومانی می‌پذیرد و
$f$ را بر آن اعمال می‌کند. به بیان دیگر، همان تابع~$f$ است:
$\lambd[x][(fx)]N$ و~$fN$ هر دو به~$fN$ کاهش می‌یابند. برای بیان این
ایده از کاهش~$\eta$ (و گسترش~$\eta$) استفاده می‌کنیم.

\begin{defn}[انقباض~$\eta$، $\eredone$]
  \ollabel{defn:beredone}
  \emph{انقباض~$\eta$} ($\eredone$) کوچک‌ترین رابطهٔ سازگار روی ترم‌هاست
  که شرط زیر را برآورده می‌کند:
  \[
    \lambd[x][M x] \eredone M \text{ مشروط بر اینکه } x \notin \FV{M}
  \]
\end{defn}

\begin{defn}[کاهش~$\beta\eta$، $\bered$] \ollabel{defn:bered}
  \emph{کاهش~$\beta\eta$} ($\bered$) کوچک‌ترین رابطهٔ بازتابی و تعدی
  روی ترم‌هاست که~$\bredone$ و~$\eredone$ را در بر می‌گیرد؛ یعنی قواعد
  بازتابی و تعدی، به‌اضافهٔ دو قاعدهٔ زیر:
  \begin{enumerate}
  \item اگر~$M \bredone N$، آنگاه~$M \bered N$. \ollabel{defn:bered3}
  \item اگر~$M \eredone N$، آنگاه~$M \bered N$. \ollabel{defn:bered4}
  \end{enumerate}
    
\end{defn}

\begin{defn}
  رابطهٔ هم‌ارزی~$\equal$ را با قاعدهٔ تبدیل~$\eta$ گسترش می‌دهیم:
  \[
  \lambd[x][f x] \equal f
  \]
  و رابطهٔ گسترش‌یافته را با~$\equal[\eta]$ نشان می‌دهیم.
\end{defn}

هم‌ارزی~$\eta$ مهم است، زیرا با برون‌گرایی ترم‌های لامبدا ارتباط دارد:

\begin{defn}[برون‌گرایی]
  رابطهٔ هم‌ارزی~$\equal$ را با قاعدهٔ~(\ext) گسترش می‌دهیم:
  \begin{center}
  اگر~$Mx \equal Nx$، آنگاه~$M \equal N$، مشروط بر اینکه~$x \notin \FV{MN}$.
  \end{center}
  و رابطهٔ گسترش‌یافته را با~$\equal[\ext]$ نشان می‌دهیم.
\end{defn}

به‌طور تقریبی، این قاعده می‌گوید دو ترم، هنگامی که به‌مثابهٔ تابع نگریسته
شوند، باید برابر دانسته شوند اگر برای آرگومان یکسان رفتار یکسانی داشته
باشند.

اکنون ثابت می‌کنیم که قاعدهٔ~$\eta$ دقیقاً برون‌گرایی را، و نه چیزی
بیشتر، فراهم می‌کند.

\begin{thm}
  $M \equal[\ext] N$ اگر و تنها اگر~$M \equal[\eta] N$.
\end{thm}

\begin{proof}
  نخست ثابت می‌کنیم~$\equal[\eta]$ تحت قاعدهٔ برون‌گرایی بسته است؛ یعنی
  قاعدهٔ~$ext$ چیزی به~$\equal[\eta]$ نمی‌افزاید. در نتیجه
  $\equal[\eta]$ رابطهٔ~$\equal[\ext]$ را در بر می‌گیرد، و اگر
  $M \equal[\ext] N$، آنگاه~$M \equal[\eta] N$.

  برای اثبات بسته‌بودن~$\equal[\eta]$ تحت~\ext، توجه کنید که برای هر
  $M \equal N$ که با قاعدهٔ~\ext{} مشتق شده باشد، مقدمهٔ
  $Mx \equal[\eta] Nx$ را داریم. سپس
  $\lambd[x][Mx] \equal[\eta] \lambd[x][Nx]$ را بنا بر یکی از
  قواعد~$\equal$ داریم؛ با اعمال~$\eta$ بر هر دو طرف،
  $M \equal[\eta] N$ را به دست می‌آوریم.

  به‌همین‌ترتیب ثابت می‌کنیم که قاعدهٔ~$\eta$ در~$\equal[\ext]$ گنجانده
  شده است. برای هر~$\lambd[x][Mx]$ و~$M$ با~$x \notin \FV{M}$، داریم
  $(\lambd[x][Mx])x \equal[\ext] Mx$، و
  $\lambd[x][Mx] \equal[\ext] M$ را از قاعدهٔ~\ext{} به دست می‌آوریم.
\end{proof}

\end{document}
